% Transcription — fonds Alexandre Grothendieck, shelfmark 31, batch 1 (pages 1-20).
% Established from the facsimile archives/batches/31/batch-01.pdf (a mirror of
% https://grothendieck.umontpellier.fr/31.pdf), per the skill
% transcribe-grothendieck.
%
% Pass: Fable 5.1 (claude-fable-5-1), 2026-09-10 — first pass, unchecked against
% the pages by a human.
%
% The batch holds three runs and four leaves of corrected typescript, each
% run opened by a cover sheet inscribed in his hand (pages 1, 9, 18), which
% become the section headings and carry no \page{}.
%
%   Pages 2, 4, 6, 8 — ink, on the backs of typescript sheets: cd_p^!(k) for
%     a field k of characteristic p, computed in the fppf topology with the
%     finite commutative p-groups, dévissé into étale, multiplicative and
%     additive pieces (Lemme 1); Lemme 2 for the multiplicative part through
%     tori and 0 → _pT → T → T → 0; the reduction to a p-extension and to μ_p
%     through Π_{k'/k}; Lemme 3, cd_p^add(k) = 0 or 1 according as k is
%     perfect or not, whose proof breaks off after « groupes α_p ».
%   Pages 10, 12, 14, 16 — ink, on the backs of the typescript leaves 30, 29,
%     28, 27, and read in that order: an « Exemple » — A a ring, f regular,
%     g regular on B = A/fA — building M = lim E_n, f-power torsion with the
%     successive torsion quotients ≃ B, but M/fM ≃ (B/gB)^{(N)}; E_{n+1} is
%     the pushout of 0 → A → A → A/fA → 0 along ξ_n ∈ E_n with ξ_n ↦ (g,0),
%     and the three inductive conditions (a_n), (b_n), (c_n) are verified.
%     The connective prose of pages 12, 14 and 16 is largely lost; the
%     formulas carry.
%   Pages 17, 15, 13, 11 — typescript leaves 27 to 30 of an exposé on
%     multiplicative groups (Corollaire 5.9 to Théorème 5.16: isotriviality,
%     constant twisted groups, the lemma on closed quasi-compact subschemes
%     of a constant twisted scheme, the theorem that a multiplicative group
%     of finite type over a locally noetherian normal base is isotrivial),
%     each struck through with a long diagonal and corrected in his hand.
%     They are transcribed in text order, which is the reverse of the
%     archive's, with the manuscript insertions in \add{} and the struck
%     words in \struck{}.
%   Page 19 — ink, fair: « Théorèmes de M. Artin », Théorème 1 on a finitely
%     generated extension K of Q (−1 a sum of squares ⟺ no ordering ⟺ not
%     embeddable in R ⟺ conditions on real points of a model ⟺ cd_2(K) < ∞)
%     and Corollaire 1 on a scheme of finite type over Z, with the bounds
%     cd(X) < 2 dim X + 1, cd(X_Q) ≤ 2 dim X_Q + 2, cd(X_R) ≤ 2 dim X. It
%     continues on page 21 (batch 2).
%
% Pages 3, 5 and 7 are typescript versos (an exposé on multiplicative
% groups, Théorème 5.1 and its corollaries) carrying nothing in his hand —
% absent. Page 20 is a blank leaf with a three-word scrap upside down at
% its foot — absent.
%
% The dating is the inventory's own, for the whole shelfmark.
\documentclass[11pt,a4paper]{article}
% Le préambule commun à toutes les transcriptions du fonds Grothendieck.
%
% Il tient en une idée : **l'incertitude se compose, elle ne se résout pas.**
% Une transcription qui lisse un mot douteux a détruit la seule chose pour
% laquelle on la faisait — un lecteur doit pouvoir savoir, à chaque endroit,
% ce qui était sur le feuillet et ce qui a été ajouté. Les six macros qui
% suivent sont donc l'appareil critique, et non de la décoration.
%
% Les mêmes commandes sont reconnues par `scripts/render.mjs`, qui produit la
% vue de lecture du site. Une macro ajoutée ici doit l'être aussi là-bas,
% faute de quoi le rendu échouera bruyamment — ce qui est voulu : un rendu
% silencieusement faux est pire qu'un rendu absent.

\ProvidesPackage{grothendieck}[2026/08/08 Transcriptions du fonds Grothendieck]

\usepackage{amsmath,amssymb,amsthm}
\usepackage{mathrsfs}
\usepackage{stmaryrd}
% Les diagrammes commutatifs sont partout dans ce fonds : c'est la langue
% même de la géométrie algébrique de Grothendieck, et une transcription qui
% les rendrait en prose ne transcrirait rien.
\usepackage{tikz-cd}
% Français par défaut — c'est la langue des feuillets — mais l'anglais est
% chargé aussi : la traduction et le résumé sont des documents anglais, et la
% typographie française (espaces avant les deux-points) y serait une faute.
% Ces documents basculent par \selectlanguage{english} en tête de corps.
\usepackage[english,french]{babel}
\usepackage{fontspec}
\usepackage{geometry}
\geometry{a4paper,margin=25mm,marginparwidth=28mm}
\usepackage{marginnote}
\usepackage{xcolor}
% ulem, not soul. `soul`'s \st and \ul are built on a character-by-character
% scan that XeLaTeX's font machinery breaks: the document fails with an
% undefined control sequence rather than degrading. ulem does the same two jobs
% and is Unicode-engine safe. `normalem` stops it hijacking \emph, which the
% transcriptions use for Grothendieck's own emphasis.
\usepackage[normalem]{ulem}
\usepackage{hyperref}
\hypersetup{colorlinks=true,linkcolor=black,urlcolor=black,pdfborder={0 0 0}}

% --- Métadonnées du lot -----------------------------------------------------
% Reprises telles quelles par le script de rendu, qui les affiche en tête de
% la vue de lecture. Elles fixent ce que le fichier prétend couvrir : sans
% elles, une transcription flotte sans point d'attache dans un fonds de seize
% mille feuillets.

\newcommand{\folder}[1]{\gdef\grothFolder{#1}}
\newcommand{\batch}[1]{\gdef\grothBatch{#1}}
\newcommand{\leaves}[2]{\gdef\grothFirst{#1}\gdef\grothLast{#2}}
\newcommand{\foldertitle}[1]{\gdef\grothFoldertitle{#1}}
\gdef\grothFolder{?}\gdef\grothBatch{?}\gdef\grothFirst{?}\gdef\grothLast{?}\gdef\grothFoldertitle{}

% --- Le résumé --------------------------------------------------------------
%
% Il ouvre la lecture modernisée, sous le titre « Résumé », et il est composé
% en petit. Ce n'est pas un ornement : c'est le seul passage du document qui
% ne soit pas des mathématiques, et un lecteur doit pouvoir le reconnaître
% avant de l'avoir lu — puis le sauter s'il connaît déjà le sujet, ou s'y
% arrêter s'il ne le connaît pas. Le filet qui le suit marque la reprise du
% corps.

\newenvironment{resume}
  {\small\setlength{\parskip}{0.45em}}
  {\par\medskip\hrule\medskip}

% --- Mots-clés ---------------------------------------------------------------
%
% En anglais, en clôture du résumé de la lecture modernisée : le vocabulaire
% moderne sous lequel chercher ce que les feuillets construisent. Le manifeste
% du site (`npm run manifest`) les extrait pour étiqueter la cote entière —
% cette ligne est la source unique des tags, il n'y a pas de fichier à côté.
\newcommand{\keywords}[1]{\par\smallskip\noindent
  {\footnotesize\textsc{Keywords} --- \textit{#1}}\par}

% --- L'appareil critique ----------------------------------------------------

% Le numéro de feuillet, en marge. C'est l'ancre qui permet de régler un
% désaccord sur une formule en regardant *un* feuillet plutôt qu'une chemise
% de six cents. Le site s'en sert aussi pour tourner les pages du fac-similé
% quand on fait défiler la transcription.
\newcommand{\leaf}[1]{%
  \marginnote{\footnotesize\color{gray!70}\textsf{#1}}%
  \label{leaf:#1}\ignorespaces}

% Illisible : on ne devine pas. Un mot inventé qui se lit comme les autres est
% le pire résultat possible.
\newcommand{\ill}{\textcolor{red!60!black}{[\,\ldots\,]}}

% Lecture douteuse : le mot est donné, et signalé comme incertain.
\newcommand{\uncertain}[1]{\uline{#1}}

% Ajout de l'éditeur — une lettre manquante, un mot sous-entendu.
\newcommand{\add}[1]{\textcolor{blue!50!black}{[#1]}}

% Biffé par Grothendieck lui-même. Ce qu'il a rayé fait partie du document :
% une rature dit souvent où la pensée a changé de direction.
\newcommand{\struck}[1]{\sout{#1}}

% Note du transcripteur, sur l'état matériel du feuillet.
\newcommand{\note}[1]{\par\smallskip{\small\color{gray!60!black}\textit{[#1]}}\par\smallskip}

% Note marginale de Grothendieck, distincte du corps du texte.
\newcommand{\marginal}[1]{\par\smallskip\noindent\hspace{1em}%
  {\small\color{gray!60!black}#1}\par\smallskip}

% --- Titre ------------------------------------------------------------------

\AtBeginDocument{%
  \begin{center}
    {\large\bfseries Fonds Alexandre Grothendieck --- cote n\textsuperscript{o}~\grothFolder}\\[2pt]
    {\itshape\grothFoldertitle}\\[4pt]
    {\small feuillets \grothFirst--\grothLast\ (lot \grothBatch)}\\[2pt]
    {\footnotesize\color{gray!60!black}Transcription établie d'après le fac-similé numérisé
     par l'Université de Montpellier.\\
     Les lectures incertaines sont soulignées, les illisibles notées [\,\ldots\,],
     les ajouts entre crochets.}
  \end{center}
  \vspace{1em}}

\endinput


\folder{31}
\batch{1}
\pages{1}{20}
\dating{[1963-1973]}
\watermark{Édition de démonstration}
\foldertitle{SGA A [SGA 4] (Résidus de rédaction) : notes manuscrites (s.d.)}

\begin{document}

\section*{$cd_p^{!}(k)$ pour $k$ de caractéristique $p$ (topologie fppf) (pages 2 à 8)}

\note{Le titre est de la main de Grothendieck, sur la chemise qui précède
(page 1) : « $cd_p^{!}(k)$ pour $k$ de car. $p$. (topolog. fppf) ». Dans ce
feuillet, « t.m. » abrège \emph{type multiplicatif}, « f.d.s. » \emph{fid.
plat} n'apparaît qu'une fois, et $cd_p^{!}$, $cd_p^{\text{ét}}$,
$cd_p^{\text{t.m.}}$, $cd_p^{\text{add}}$ désignent les dimensions
cohomologiques calculées avec les $p$-groupes finis de chacun des trois
types dévissés à la page 2.}

\page{2}
$k$ corps de car $p > 0$

On désigne par $cd_p(k)$ la dimension cohomologique, calculée pour les
faisceaux de $p$-torsion pour la topologie finie \struck{étale} fid.\ plate
sur $k$. Regardons d'abord le $cd_p^{!}(k)$, calculée avec les groupes
algébriques (commutatifs) finis sur $k$, d'ordre une puiss.\ de $p$.
\uncertain{Tout} tel groupe algébrique se dévisse en
\[
G' \subset G'' \subset G
\]
où $G'' = G^{\circ}$ donc $G/G''$ est un groupe étale, et où $G'$ est
(radiciel) de t.m., tandis que $G''/G'$ (radiciel) \add{est} de type
\uncertain{additif}. Donc

\textbf{Lemme 1} : $cd_p^{!}(k) \leqslant n \Longleftrightarrow H^i(k,G) = 0$
pour $i > n$, lorsque $G$ est un $p$-groupe fini sur $k$ de l'un des types
suivants :
\begin{enumerate}
\item[a)] $G$ étale sur $k$, [i.e.\ $\Longrightarrow cd_p^{\text{ét}}(k) \leqslant n$]
\item[b)] $G$ radiciel de t.m.\ sur $k$
\item[c)] $G$ \add{radiciel} de type additif.
\end{enumerate}
\note{Dans a), le mot « $G$ » est repassé ; dans c), « radiciel » est écrit
par-dessus un mot biffé.}

\page{4}
Notons que $cd_p^{\text{ét}}(k)$ est égal à $0$ ou $1$, suivant que
\struck{toute} extension \uncertain{finie} sép.\ de $k$ est
\uncertain{d'ordre} premier à $p$, ou non.
\note{Le mot « finie » est cerclé et rattaché par un trait à la ligne, comme
une insertion ; le mot biffé avant « extension » est douteux.}

Si $G$ est de t.m., donc $G$ est dual d'un groupe fini étale, ce dernier
est quotient d'un \uncertain{module galoisien} \uncertain{libre} de type fini
sur $\mathbb{Z}$, donc $G$ est contenu dans un tore, donc on a une suite
exacte
\[
0 \longrightarrow G \longrightarrow T \longrightarrow T' \longrightarrow 0
\]
$T$ et $T'$ des tores. D'où le

\textbf{Lemme 2} : Pour que $cd_p^{\text{t.m.}}(k)$ soit $\leqslant n$, il
f.\ et s.\ que les conditions suivantes soient vérifiées : pour tt tore $T$
sur $k$, $H^i(k,T)(p)$ est nul pour $i > n$, et $p$-divisible si $i = n$.
[Ou encore : pour tt extension finie sép.\ $k'$ de $k$, \ill{}
\struck{$\mathrm{Br}(k)(p)$} $H^i(k', G_m) = 0$ si $i > n$, et
$H^i(k', G_m)$ est $p$-divisible si $i = n$.]
\note{La fin du crochet est écrite dans la marge gauche, de bas en haut.}
Comme on a
\[
0 \longrightarrow {}_{p}T \longrightarrow T \xrightarrow{\;p\;} T \longrightarrow 0
\]
d'où \struck{$H^i({}_pT) \equiv \mathrm{Im}\, H^i({}_pT)$}
\note{Formule biffée de deux traits ; la lecture en est douteuse.}

\page{6}
\[
H^i({}_pT) \longrightarrow {}_pH^i(T) \longrightarrow 0
\]
\[
0 \longrightarrow H^i(T)_p \longrightarrow H^{i+1}({}_pT)
\]
on voit que ces conditions sont nécessaires. Pour la suffisance, on tire
d'abord le fait que $H^{n+i}(G)$ \struck{\ill{}} $= 0$ pour $i \geqslant$
\uncertain{$n+2$}, reste à vérifier $H^{n+1}(G) = 0$ donc que
$H^n(T) \longrightarrow H^n(T') \longrightarrow 0$ surjectif.
\note{Les deux termes $H^n(T)$, $H^n(T')$ portent un indice griffonné puis
biffé, illisible.} \struck{\ill{}}
\note{Suit un bloc de cinq lignes entièrement barré de traits croisés, dont
seuls se lisent « $p$-torsion », « $H^n(T)(p) \to H^n(T')(p)$ » et, sur la
dernière ligne, « $p$-divisible » ; le reste est illisible sous les
ratures.}

Notons que si $k'/k$ une extension finie séparable de $k$, alors pour tt
$G'$ sur $k'$, on a
\[
H^i(k, \textstyle\prod_{k'/k} G') \simeq H^i(k, G),
\]
\note{Sic pour le membre de droite ; le sens demande $H^i(k', G')$.}
ce qui prouve que l'hypothèse faite sur les tores sur $k$ est
\uncertain{stable} par extension de base séparable. Soit $k'$ une telle
extension de degré premier à $p$, alors l'homomorphisme \uncertain{naturel}
\[
\textstyle\prod_{k'/k} G_{k'} \longrightarrow G
\]
est \struck{surjectif} un épim.

\page{8}
[\textbf{N.B.} $\prod_{k'/k} G_{k'}$ est de t.m.] donc \ill{} :
$H^{n+1} = 0$ sur les groupes de t.m.\ \ill{} on voit que
\[
H^{n+1}(\textstyle\prod_{k'/k} G_{k'}) \longrightarrow H^{n+1}(G) \longrightarrow 0
\quad\text{surjectif,}
\]
\[
\simeq H^{n+1}(G_{k'})
\]
\note{Le second membre est écrit sous le premier terme de la ligne
précédente, avec un signe $\simeq$ qui le relie à
$H^{n+1}(\prod_{k'/k} G_{k'})$.}
donc pour vérifier que $H^{n+1}(G) = 0$, il suffit de le faire pour
$H^{n+1}(G_{k'})$. On nous ramène au cas où $G$ est « splitté » par une
$p$-extension, ce qui implique que $G$ se dévisse en groupes isomorphes à
$\mu_p$. Donc on est ramené à vérifier que $H^{n+1}(\mu_p) = 0$, ce qui se
voit sur la suite exacte
\[
\cdot \longrightarrow \mu_p \longrightarrow G_m \longrightarrow G_m \longrightarrow 0
\]

\textbf{Lemme 3} \struck{\ill{}} On a
\[
cd_p^{\text{add}}(k) =
\begin{cases}
0 & \text{si } k \text{ parfait} \\
1 & \text{si } k \text{ imparfait}
\end{cases}
\]
\note{Au-dessus de l'exposant « add » est écrit « rad », d'une plume plus
fine.}

En effet on \uncertain{peut} \ill{} radiciel additif se dévisse en groupes
$\alpha_p$ \ldots
\note{Le feuillet s'arrête sur cette ligne inachevée.}

\section*{Généralités sur les produits \ill{} : un module $M$ de $f$-torsion dont $M/fM$ n'est pas de type fini (pages 10 à 16)}

\note{Le titre est de la main de Grothendieck, sur la chemise qui précède
(page 9) : « Généralités sur les produits \ill{} » ; le dernier mot, de
lecture douteuse, se termine en « -adiques ». Le sous-titre est de
l'éditeur. Ces quatre feuillets sont écrits au dos des feuillets de
tapuscrit des pages 11, 13, 15 et 17. Ici ${}_{f}M$, ${}_{f^n}M$
désignent, comme sur les feuillets, les sous-modules de $M$ annulés par
$f$, par $f^n$ : le $f$ est écrit en indice à gauche du module. $B$ désigne
$A/fA$ tout au long.}

\page{10}
\textbf{Exemple.} $A$ anneau, $f \in A$ élément $A$-régulier, $g \in A$
élément $(A/fA)$-régulier, on suppose $A/fA/gA \neq 0$ i.e.\ $B/gB \neq 0$.
\note{« $= B$ » est écrit au-dessus de $(A/fA)$.}
[\ill{} $A$ noethérien régulier de dim $2$, $A/fA$ noeth.\ régulier de dim
$1$, $A/fA/gA$ noeth.\ régulier de dim $0$, p.\ ex.\ $A$ \uncertain{local
régulier} de dim $2$, $(f,g)$ \uncertain{système de paramètres}]

On va construire un module $M$ sur $A$ tel que
\begin{enumerate}
\item[(i)] $M = \bigcup_{n \geqslant 0} {}_{f^n}M$, \quad
${}_{f^{n+1}}M / {}_{f^n}M \simeq A/fA$ \quad [en particulier, $M$ est de
\uncertain{support} $\subset V(f)$ \ill{} $fA$-\uncertain{cofini}]
\item[(ii)] $M/fM \simeq (A/fA + gA)^{(\mathbb{N})} = (B/gB)^{(\mathbb{N})}$,
donc $M/fM$ n'est pas $fA$-\uncertain{cofini}.
\end{enumerate}
\note{Dans (ii), « $A/fA + gA$ » est écrit tel quel, pour $A/(fA+gA)$.}

\textbf{Construction.} On construit par récurrence les $E_n$, l'hypothèse
de récurrence étant :
\note{À gauche de ces deux lignes, un \ill{} de deux mots, dont le
premier semble être $fM$.}
\begin{enumerate}
\item[(1)] $f^n E_n = 0$
\item[(2)] ${}_{f^{k+1}}E_n / {}_{f^k}E_n \simeq A/fA = B \quad (0 \leqslant k \leqslant n-1)$
\item[(3)] $E_n / fE_n \simeq$ \struck{$A/fA$} $B + (B/gB)^{n-1}$
\end{enumerate}
\note{Les trois conditions sont accolées ; la page 14 les appelle $(a_n)$,
$(b_n)$, $(c_n)$.}

\page{14}
On construit $E_{n+1}$ comme extension de $A/fA$ par $E_n$, déduite de
l'extension \uncertain{standard} de $A/fA$ par $A$ \ill{} par
l'homomorphisme \struck{\ill{}} $A \to E_n$ défini par un \ill{} $\xi_n$ de
$E_n$ (l'extension $E_{n+1}$ ne dépend que de $\xi_n$ \uncertain{mod}
$fE_n$).

\begin{tikzcd}[column sep=small]
0 \arrow[r] & A \arrow[r, "f"] \arrow[d, "\xi_n"'] & A \arrow[r] \arrow[d] & A/fA \arrow[r] \arrow[d, no head, "\|" description] & 0 \\
0 \arrow[r] & E_n \arrow[r] & E_{n+1} \arrow[r] & A/fA \arrow[r] & 0
\end{tikzcd}

\note{Le $0$ de gauche est, sur les deux lignes, un simple point.}

Notons que ${}_fE_n \subset {}_{f^{n-1}}E_n$ donc
\[
0 \longrightarrow {}_{f^{n-1}}E_n / {}_fE_n \longrightarrow E_n / {}_fE_n
\longrightarrow E_n / {}_{f^{n-1}}E_n \longrightarrow 0
\]
\[
E_n / {}_{f^{n-1}}E_n \simeq A/fA = B
\]
\note{L'isomorphisme est écrit sous le dernier terme de la suite, relié
par un $\wr$.}

donc \ill{} des $B$ \struck{et} \struck{\ill{}} ($B$-\ill{}) \ill{} (ii)
$B$ \ill{}, \ill{} de \ill{} $2$. Ce \struck{\ill{}} \ill{} $\lambda$ pour
(\ill{}), et s'il \ill{} $\lambda$ \ill{}. Donc l'iso\ill{} \ill{}
[i.e.\ ${}_{f^{n-1}}E_n / {}_fE_n \simeq (B/gB)^{n-1}$] d'extension
\ill{}.
\note{Six lignes de prose serrée, dont seuls les mots donnés se lisent ;
la formule entre crochets est nette.}

Nous choisirons $\xi_n =$ \struck{\ill{}} $(g, 0)$. \struck{Alors}

Vérifions les propriétés $(a_{n+1})$, $(b_{n+1})$, $(c_{n+1})$.

Le premier résulte \ill{} $(a_n)$ et $f \cdot A/fA = 0$
\note{Le feuillet s'arrête ici ; la vérification se poursuit page 12.}

\page{12}
la deuxième \struck{\ill{}} \ill{} voulue\uncertain{s} :
\[
E_n = {}_{f^n}E_{n+1},
\]
\note{Le signe est un $=$ barré d'un trait oblique ; l'argument qui suit
établit l'égalité.}
\ill{} $\supset$, donc : \ill{} Soit $x \in E_n$, $h \in A$, \struck{\ill{}}
tels qu'il existe $\lambda \in A$, avec
\[
f(x,h) = \lambda(\xi_n, f)
\]
\ill{}, \ill{} $h \in fA$. \ill{}, \ill{}.
$f^n h = \lambda f$ \ill{} $\lambda = f^{n-1}h$, (car $f$ est $A$-régulier)


$f^n x = \lambda \xi_n$ \struck{\ill{}} \ill{} i.e.


$0 = f^{n-1} h \xi_n$, \ill{} $(f^{n-1}\xi_n) \in {}_fE_n$ \ill{},

\uncertain{d'où} \ill{} ${}_fE_n = A/fA = B$ \ill{} $h$), donc \ill{}
\ill{} \ill{} $\xi_n \notin {}_{f^{n-1}}E_n$ i.e.\ \struck{$f^{n}$}
$y_n \neq 0$, donc $h y_n = 0 \Longrightarrow h \in fA$, cqfd.
\note{$y_n$ n'est pas défini sur le feuillet ; le contexte demande
$y_n = f^{n-1}\xi_n$.}

\uncertain{Enfin}, \uncertain{vérifions} \ill{} :
\begin{align*}
E_{n+1}/fE_{n+1} &\simeq (E_n \times A) \big/ \big( f(E_n \times A) + A(\xi_n, f) \big) \\
&\simeq (E_n/fE_n \times A/fA) \big/ A(\bar{\xi}_n, f) \\
&\simeq \big( [B \times (B/gB)^{n-1}] \times B \big) \big/ B(g, 0, 0) \\
&\simeq B/gB \times (B/gB)^{n-1} \times B \simeq B \times (B/gB)^{n}.
\end{align*}

\page{16}
De plus la construction \ill{} donne un \uncertain{homomorphisme}
\[
E_n \xrightarrow{\;u_n\;} E_{n+1}
\]
\uncertain{d'où} par \ill{} un \ill{}
\[
u_n \otimes_A B : E_n/fE_n \longrightarrow E_{n+1}/fE_{n+1}.
\]
Soit \emph{\ill{}} \ill{} $g(E_n/fE_n)$, \uncertain{et} \ill{} \ill{}
\uncertain{identifié} dans $($\struck{\ill{}} $E_{n+1}/fE_{n+1} \otimes_B
B/gB)$ \ill{} \struck{\ill{}} \ill{} $Q_n \otimes B/gB$, \ill{} \ill{}
(\ill{}).
\note{Deux lignes entières sont biffées d'un long trait ; $Q_n$ n'est
défini que par la ligne suivante.}

Donc
\[
\varinjlim Q_n = \varinjlim\, Q_n \otimes_B B/gB \simeq \varinjlim\, (B/gB)^{n-1}
\simeq (B/gB)^{(\mathbb{N})}.
\]

Prenons $M = \varinjlim E_n$

\struck{\ill{}} $M$ \uncertain{satisfait} aux conditions voulues.

\section*{Quatre feuillets de tapuscrit corrigé : groupes de type multiplicatif isotriviaux, 5.9 à 5.16 (pages 17, 15, 13 et 11)}

\note{Feuillets numérotés 27 à 30 par le dactylographe, rangés à rebours
dans le dossier : la page 17 porte le feuillet 27, la page 15 le feuillet
28, la page 13 le feuillet 29 et la page 11 le feuillet 30. Ils sont donnés
ici dans l'ordre du texte, qui se poursuit d'un feuillet à l'autre. Chacun
est barré d'un long trait diagonal. Les mots surfrappés par la machine et
les mots biffés sont donnés barrés ; les ajouts de la main de Grothendieck
sont donnés entre crochets, comme le sont ailleurs les ajouts de
l'éditeur. Les numéros des énoncés ont été corrigés à la main
sur les quatre feuillets, le chiffre des unités étant repassé (5.8 en 5.9,
5.9 en 5.10, etc.) ; on donne les numéros corrigés. Les versos de ces
feuillets portent les notes manuscrites des pages 10, 12, 14 et 16.}

\page{17}
\textbf{Corollaire 5.9.} Les foncteurs 5.7.\ induisent des anti-équivalences
quasi-inverses l'une de l'autre entre la catégorie des groupes de type
multiplicatif de type fini \add{$H$} sur $S$, et la catégorie des groupes
constants tordus à engendrement fini \add{$R$} sur $S$. Tout tel groupe $R$
est quasi-isotrivial.

\textbf{Corollaire 5.10.} Soient $H, G$ deux \struck{été} $S$-préschémas en
groupes de type multiplicatif et de type fini. Alors
$\underline{\mathrm{Isom}}_{S\text{-gr}}(H,G)$ est représentable par un
sous-préschéma ouvert \add{et fermé} de
$\underline{\mathrm{Hom}}_{S\text{-gr}}(H,G)$, et c'est un \struck{schéma}
$S$-préschéma constant tordu. En particulier,
$\underline{\mathrm{Aut}}_{S\text{-gr}}(H)$ est représentable et est un
groupe constant tordu (en général non commutatif) sur $S$.

Cela résulte de 5.8.\ et de Exp \add{VIII} 1.6.
\note{Le numéro de l'exposé est repassé à l'encre par-dessus le numéro
typographié.}

\textbf{Proposition 5.11.} Soient $S$ un préschéma, $R$ un groupe constant
tordu sur $S$, $H = D_S(R)$ le groupe de type multiplicatif qu'il définit.
Considérons les conditions suivantes :
\begin{enumerate}
\item[(i)] $H$ est isotrivial, (i.e.\ $R$ est isotrivial).
\item[(ii)] $R$ est \struck{somme disjointe de} \add{réunion de}
sous-préschémas $R_i$, qui sont \add{à la fois ouverts et fermés,}
\struck{quasi-compacts} (et alors nécessairement finis) sur $S$.
\item[(iii)] Les composantes connexes de $R$ sont finies sur $S$.
\end{enumerate}
Alors (i) $\Longrightarrow$ (ii) $\Longrightarrow$ (iii), on a \add{(i)
$\Longleftrightarrow$} (ii) $\Longleftrightarrow$ (iii) si $S$ est
localement noethérien, enfin (i) $\Longleftrightarrow$ (ii) si $R$ est
\struck{de} à engendrement fini i.e.\ si $H$ est de type fini\add{, du
moins si $S$ est quasi-compact ou si ses composantes connexes sont
ouvertes.}
\note{L'ajout « (i) $\Longleftrightarrow$ » est écrit au-dessus de la ligne,
avec un trait qui le rattache à l'endroit qui précède « (ii)
$\Longleftrightarrow$ (iii) ».}

\struck{Démonstration. (i) $\Longrightarrow$ (ii) Soit $S'\to S$ fini étale
surjectif qui splitte $H$ donc $R$, de sorte que $R'$ est somme}
Décomposant d'abord $S$ en préschéma somme de préschémas $S_i$ sur chacun
desquels $H$ est de type constant (cf Exp \add{IX} 1), on est ramené au
cas où $H$ \struck{est donc} donc $R$ est de type constant. Nous aurons
besoin \struck{\ill{}} d'un

\struck{Lemme 5.11.}

\page{15}
\textbf{Lemme 5.12.} Soient $S$ un préschéma, $R$ un préschéma constant tordu
sur $S$. Alors tout sous-préschéma fermé \struck{\ill{}} \add{$Z$} de $R$
qui est \struck{de type fini} \add{quasi-compact} sur $S$, est fini sur $S$.

En effet, on est ramené au cas où $R$ est constant, donc de la forme $I_S$,
où $I$ est un ensemble, donc réunion filtrante croissante des $J_S$, où $J$
parcourt les parties finies de $I$. On peut supposer de plus $S$ affine,
alors \struck{\ill{}} $Z$ est quasi-compact, donc contenu dans un des $J_S$.
Comme $J_S$ est fini sur $S$, il en est de même de $Z$.

Le lemme 5.12.\ établit déjà l'assertion entre parenthèses dans 5.11.\ (ii).
\struck{L'équivalence (ii) et (iii) est} \add{L'implication (ii)
$\Longrightarrow$ (iii), et l'implication inverse} dans le cas $S$
localement noethérien est alors claire. Prouvons que (i) $\Longrightarrow$
(ii), pour ceci, soit $S' \to S$ un morphisme étale fini surjectif qui
splitte $H$ donc $R$, de sorte que $R'$ s'écrit \struck{comme}
$M_{S'} =$ \struck{\ill{}} $\coprod_{m \in M} R'_m$, où \struck{tout} les
$R'_m$ sont des ouverts disjoints de $R'$, isomorphes à $S'$. Soit
$g : R' \to R$ la projection, qui est un morphisme fini étale surjectif,
\add{donc un morphisme ouvert et fermé ;} alors les $R_m = g(R'_m)$ sont
des parties à la fois ouvertes et fermées de $R$, et évidemment
quasi-compacts sur $S$ puisque les $R_m$ le sont. \struck{Enfin, prouvons
(ii) $\Longrightarrow$ (i)} \add{$S$ quasi compact ou connexe.}

Enfin, prouvons (ii) $\Longrightarrow$ (i) lorsque $H$ est de type fini sur
$S$. Soit $M$ son type, on peut écrire $M =$ \struck{\ill{}} $\mathbb{Z}^r
\times N$, où $r$ est un entier $\geqslant 0$ et $N$ un groupe fini.
Considérons le préschéma
\[
P = \underline{\mathrm{Isom}}_{S\text{-gr}}(H,G) \subset
\underline{\mathrm{Hom}}_{S\text{-gr}}(H,G) = Q
\]
où $G = D_S(M)$, cf 5.8.\ et 5.10. \struck{Alors $Q$ est isomorphe} On a un
isomorphisme
\[
Q \simeq \underline{\mathrm{Hom}}_{S\text{-gr}}(G, \underline{G}_m^{\,r})
\times \underline{\mathrm{Hom}}_{S\text{-gr}}(G,N) \simeq R^r \times E,
\]
où $E = \underline{\mathrm{Hom}}_{S\text{-gr}}(G,N)$ est fini sur $S$. Il
s'ensuit que $Q$ est \struck{somme de} \add{réunion de sous-}préschémas à la
fois ouverts et fermés $Q_i$ finis sur $S$. \struck{De même, Il s'ensuit que}
\add{Donc} $P$ est réunion \struck{disj} de \struck{prés} sous-préschémas à
la fois ouverts et fermés $P_i$ \add{$= P \cap Q_i$} finis sur $S$. Comme ils
sont étales sur $S$, leurs

\page{13}
images dans $S$ sont des parties $S_i$ à la fois ouvertes et fermées, et
elles recouvrent $S$. Si $S$ est connexe ou quasi-compact, il existe donc un
ensemble fini d'indices \struck{\ill{}} $i$ tels que les $S_i$ recouvrent
$S$, soit $S'$ la réunion des $P_i$ correspondants, \struck{qui est donc tel
que} \add{Alors} $S' \to S$ \struck{est} \add{est} fini étale surjectif, et
posant $P' = P \times_S S' = \underline{\mathrm{Isom}}_{S'\text{-gr}}(H',G')$,
on voit que $P'$ a une section sur $S'$, donc il existe un isomorphisme de
$S'$-groupes $H' = H \times_S S' \xrightarrow{\;\sim\;} G' = G \times_S S' =
D_{S'}(M)$, ce qui prouve que $S'$ splitte $H$. Le cas où les composantes
\struck{irréductibles} \add{connexes} de $S$ sont \struck{\ill{}} ouvertes se
ramène aussitôt au cas où $S$ est connexe. --- Remarquons d'ailleurs qu'on
peut, lorsque $H$ est de type fini sur $S$ \struck{\ill{}} et de type
constant, donner le critère d'isotrivialité suivant (où il n'est plus
nécessaire de faire aucune restriction sur $S$) : $R$ est réunion d'une
\emph{suite} de parties à la fois ouvertes et fermées finies sur $S$.

\textbf{Lemme 5.13.} Soient $S$ un préschéma localement noethérien et
connexe, \add{$P$} un $S$-préschéma constant tordu quasi-isotrivial, $Z$ une
partie à la fois ouverte et fermée de \add{$P$}, telle qu'il existe un
$s \in S$ tel que \add{$Z_s$} soit fini. Alors \add{$Z$} est fini sur $S$.
\note{Dans cet énoncé les lettres typographiées sont repassées à l'encre :
$R$ corrigé en $P$ aux deux endroits, et $R_s$, $R$ corrigés en $Z_s$, $Z$ ;
la marge gauche porte « $Z_s$ » et « $Z$ » en rappel.}

\add{Ne supposons pas d'abord $S$ connexe :} \struck{\ill{} $S$ est supposé
connexe, et} soit $U$ l'ensemble des $s \in S$ tels que $Z_s$ soit fini,
nous allons prouver que $U$ est à la fois ouvert et fermé, et que $Z|U$ est
fini sur $U$. C'est là évidemment un \struck{\ill{}} énoncé essentiellement
équivalent à 5.13.\ mais qui a l'avantage d'être « de nature locale » sur
$S$ pour la topologie étale (disons), ce qui nous ramène au cas où \add{$P$}
est constant, i.e.\ de la forme $I_S$, où $I$ est un ensemble. (NB
l'hypothèse localement noethérienne n'est pas perdue par un changement de
base étale ; c'est là où on utilise la quasi-isotrivialité de \add{$P$} sur
$S$). On peut de plus supposer $S$ connexe, \struck{Mais alors}

\page{11}
\struck{\ill{}} puisque ses composantes connexes sont ouvertes ($S$ étant
localement noethérien). Mais alors on a nécessairement $Z = J_S$, où $J$ est
une partie de $I$, et on a \struck{\ill{}} donc $U = \emptyset$ ou $U = S$,
suivant que $J$ est infini ou fini, ce qui donne la conclusion voulue.

\textbf{Corollaire 5.14.} Soi\add{ent} $S$ un préschéma localement
noethérien et \emph{normal}, \add{$P$} un $S$-préschéma \struck{localement}
constant tordu quasi-isotrivial. Alors les composantes connexes de
\add{$P$} sont finies sur $S$.

Comme \add{$P$} est étale sur $S$, on sait que \add{$P$} est normal (SGA I
9.5.) donc ses composantes \struck{irréductibles} connexes sont ses
composantes irréductibles. On peut évidemment supposer $S$ connexe donc
irréductible, soit $s$ son point générique. Comme $R$ est plat sur $S$,
ses points maximaux, i.e.\ les points génériques $x_i$ des composantes
irréductibles $R_i$ de $R$, sont les points maximaux de $R_s$, et la fibre
\struck{dans} générique $R_{i\,s}$ de $R_i$ est la composante irréductible
de $R_s$, adhérence de $x_i$ dans $R_s$. Or $R_s$ étant étale sur $k(s)$,
est somme disjointe de spectres d'extensions finies séparables de $k(s)$,
donc \struck{les} pour tout $i$, $R_{i\,s}$ est fini (en fait, réduit à un
seul point). Donc 5.13.\ s'applique et montre que les $R_i$ sont finis sur
$S$.
\note{Dans l'énoncé de 5.14 le $R$ typographié est corrigé en $P$, mais la
démonstration garde $R$ ; on transcrit tel quel.}

\textbf{Remarque 5.15.} Le raisonnement s'applique, plus généralement, si on
suppose seulement $S$ géométriquement unibranche au lieu de normal, i.e.\
que le morphisme canonique $\widetilde{S} \to S$, où $\widetilde{S}$ est le
normalisé de $S_{\text{réd}}$, est radiciel. Donc le théorème 5.16.\
ci-dessous est valable également sous cette condition plus générale.

\textbf{Théorème 5.16.} Soit $S$ un préschéma localement noethérien normal.
Alors tout groupe de type multiplicatif \struck{quasi-isotrivial} \add{et de
type fini} sur $S$ \struck{\ill{}} est \emph{isotrivial}.

On peut supposer en effet \add{$S$ connexe, donc} $H$ de type constant $M$.
\add{[Il suffit d'appliquer 5.13.\ avec $P =
\underline{\mathrm{Isom}}_{S\text{-gr}}(H,G)$, où $G = D_S(M)$, en
raisonnant comme dans 5.11.\ (ii) $\Longrightarrow$ (i) \struck{\ill{}}, en
utilisant 5.14 pour $R = D_S(H)$.]}
\note{Le crochet et ses deux lignes manuscrites sont au bas du feuillet,
la seconde entièrement de la main de Grothendieck ; en marge, un « M »
cerclé et « 5.11. ».}

\section*{Théorèmes de M. Artin sur cd (page 19)}

\note{Le titre est de la main de Grothendieck, sur la chemise qui précède
(page 18) : « Théorème de M. Artin sur cd ». La page 19 porte en tête, de
sa main et souligné, « Théorèmes de M. Artin ». Le feuillet suivant du
dossier (page 21, batch 2) poursuit avec les Corollaires 2 et 3.}

\page{19}
\marginal{\ill{} ?!}
\note{Note marginale de quatre lignes, écrite en biais dans la marge gauche
à hauteur du Théorème 1, qui se termine par « ?! » ; le reste est
illisible.}

\textbf{Th.\ 1} Soit $K$ une extension de type fini de $\mathbb{Q}$,
Conditions équivalentes :
\begin{enumerate}
\item[(i)] $-1$ est une somme de carrés dans $K$
\item[(i bis)] $K$ n'admet pas de structure de corps ordonné
\item[(i ter)] $K$ n'est pas isomorphe à un sous-corps de $\mathbb{R}$
\item[(ii)] \struck{Il existe un modèle} $X$ \struck{de $K$ qui} n'a pas de
pt simple réel.
\end{enumerate}
\note{Sur la ligne (ii), « $X$ un modèle » est écrit au-dessus, relié par
un trait courbe à la fin de la ligne ; suit un bloc de trois lignes
encadré et barré de deux traits obliques :}
\begin{quote}
\struck{(ii bis) Tout modèle $X$ de $K$ est sans point simple réel.
(ii ter) Il existe un modèle [propre] \struck{sans point simple}
\uncertain{simple}, sans pt réel. (ii quater) Tout modèle (propre)
\uncertain{simple} est sans pts réels.}
\end{quote}
\begin{enumerate}
\item[(ii bis)] \struck{(si $X$ est un modèle)} $\dim X(\mathbb{R}) < n =
\deg.\ \text{tr.}\ K/\mathbb{Q}$
\note{« topologique » est cerclé et inséré au-dessus de $\dim$.}
\item[(ii ter)] $X(\mathbb{R})$ \struck{contenu dans} \uncertain{rare}
dans $X_{\mathbb{C}}$ \struck{dans une partie} \ill{} top.\ de Zariski.
\item[(ii quater)] (si $X$ simple $/\mathbb{Q}$) $X(\mathbb{R}) = \emptyset$
\item[(iii)] $cd_2(K) < +\infty$
\item[(iii bis)] (si $X_{\mathbb{R}}$ simple sur $\mathbb{R}$)
\note{L'indice de $X$ est douteux.}
$cd_2(X) < +\infty$
\end{enumerate}
\note{Sous « (iii bis) », une seconde étiquette est grattée jusqu'à être
illisible.}

\textbf{Corollaire 1} Soit $X$ un préschéma \struck{resp.\ un} de type fini
sur $\mathbb{Z}$ \struck{(resp.\ sur $\mathbb{Q}$)} Conditions équivalentes :
\begin{enumerate}
\item[(i)] $X(\mathbb{R}) = \emptyset$
\item[(ii)] $cd_2(X) < +\infty$
\item[(ii bis)] \struck{$cd_2(X) \leqslant 2\dim X + 1$ (resp.\
$\leqslant 2\dim X + 2$)}
\note{L'étiquette « (ii bis) » est biffée avec la ligne.}
\item[(iii)] $cd_2(X_{\mathbb{Q}}) < +\infty$
\end{enumerate}
\struck{S'il en est ainsi, on a}
\begin{enumerate}
\item[(iv)] $cd_2(X_{\mathbb{R}}) < +\infty$
\end{enumerate}
Alors
\[
cd(X) < 2\dim X + 1, \qquad cd(X_{\mathbb{Q}}) \leqslant 2\dim X_{\mathbb{Q}} + 2,
\qquad cd(X_{\mathbb{R}}) \leqslant 2\dim X
\]
\note{Les deux dernières inégalités sont écrites en bas de page, l'une sous
l'autre ; le premier signe est un $<$, les deux autres des $\leqslant$.
L'énoncé se poursuit page 21.}

\end{document}
